% ============================================================
%  Experiments
% ============================================================
\section{Experiments}
\label{sec:exp}

\subsection{Experimental Setup}
All methods share one frozen CLIP ViT-B/16 encoder unless stated, so rows of
Table~\ref{tab:main} differ in post-processing and auxiliary inputs, not in the
representation; Section~\ref{sec:cross} repeats the comparison on two further frozen
encoders (ViT-L/14 and SigLIP2). Heads are optimised with AdamW for $1{,}500$ iterations at batch size
$256$, learning rate $10^{-3}$, weight decay $10^{-4}$ and the Huber loss with its
default transition; three seeds share a fixed training subset, row order and therefore
$\mu,\sigma$, and hashes of the manifest, feature cache and row order are verified to
agree. The primary metric is AUROC over pairs; AUPRC is secondary; 2AFC is reported
separately because it is not interchangeable with AUROC. Values are mean $\pm$
standard deviation over three seeds; the data subset is fixed. Comparisons use a cycle-clustered paired
bootstrap with Holm correction over the $24$ cells;
CSLS's $\lambda$ is selected on the same validation split; re-selecting it on a
$0$--$30$ grid leaves the Quantity/Entity ordering unchanged. The Action facet's
official validation split contains two cycles, so for that facet only $20\%$ of the
training cycles ($80$ of $399$) are held out, with no image shared with the retained
cycles; this is a single fixed draw and a protocol deviation. Its $166$ test cycles
are untouched, and the Action cells are marked $\dagger$ and treated as exploratory.

\subsection{Comparison with Same-Backbone Detectors}
Table~\ref{tab:main} compares CRD with a cosine threshold, CSLS, a GMM density
detector, LEMoN-style and TRACED-style reproductions, two parameter-free
feature-debiasing baselines (mean-centred and all-but-the-top (ABTT) cosine), and a
supervised probe used as a reference ceiling. The
two debiasing baselines are the null hypothesis: they ask whether the gain comes from
the learned heads or merely from removing first-order feature structure. Both are
fitted on the training split only, with no gradient training; the single
hyper-parameter ($k$ for ABTT) is selected on the same validation split as $\lambda$.
LEMoN uses the official adapter on the same cached CLIP features and test split:
IndexFlatIP neighbour retrieval with $k{=}5$, the official
$d_1+\beta d_n+\gamma d_m$ score, and hyperparameters fitted on the validation split.
TRACED is a mechanism-level reproduction: it builds a deletion-only caption trajectory
of three steps; because its token-replacement operator is not reproduced, it is treated
as a lower bound.

\begin{table*}[t]
\centering
\footnotesize
\setlength{\tabcolsep}{9pt}
\caption{AUROC on FGBench v1-strict, frozen CLIP ViT-B/16. CRD is mean $\pm$ std over
seeds; baselines are single-run. $^\dagger$Action uses an internal validation split.
GMM/mean-centred/ABTT are controls; the LR probe is reported in the text. The Holm
family is five method rows $\times$ five facets minus GMM/Action ($24$ cells; four
GMM duplicates leave $20$ independent).}
\label{tab:main}
\begin{tabular}{lccccc}
\toprule
Method & Entity & Quantity & Attribute & Action$^\dagger$ & Relation \\
\midrule
CLIP cosine                & 0.9991 & 0.8305 & 0.6806 & 0.5913 & 0.5413 \\
CLIP + CSLS                & \textbf{0.9996} & 0.8882 & 0.6866 & 0.5869 & 0.5417 \\
LEMoN (official)           & 0.9990 & 0.8266 & 0.6755 & 0.5923 & 0.5218 \\
TRACED-style               & 0.9985 & 0.8412 & 0.6875 & 0.6102 & 0.5356 \\
\midrule
\multicolumn{6}{l}{\emph{Controls (not detectors)}}\\
GMM (monotone)             & 0.9991 & 0.8305 & 0.6806 & --- & 0.5413 \\
Mean-centred cosine        & 0.9989 & 0.8570 & 0.6790 & 0.5890 & 0.5389 \\
ABTT cosine                & 0.9985 & 0.8481 & 0.6765 & 0.5394 & 0.5464 \\
\midrule
\textbf{CRD (ours)}        & \shortstack{0.9995\\\scriptsize$\pm$0.0000} & \shortstack{0.8837\\\scriptsize$\pm$0.0011} & \shortstack{0.7187\\\scriptsize$\pm$0.0001} & \shortstack{0.6145\\\scriptsize$\pm$0.0009} & \shortstack{\textbf{0.5530}\\\scriptsize$\pm$0.0001} \\
\bottomrule
\end{tabular}
\end{table*}

On Attribute, CRD improves AUROC by $3.81$ points on the seed mean (arithmetic
difference; the bootstrap mean is $3.84$) and by $3.12$ to $4.32$ points over the
five detector/control rows in Table~\ref{tab:main} (cosine, CSLS, GMM, LEMoN official
and TRACED-style), with a 95\% CI of $[3.73,3.96]$ against cosine. The absolute gain is larger on Quantity ($+5.32$), but CSLS already
accounts for most of it. On Action CRD reaches $0.6145$ against $0.6102$ for TRACED,
a paired-bootstrap difference of $+0.43$\,pp (95\% CI $[-1.75,2.54]$), which does
not survive Holm; we claim parity rather than superiority. On Quantity and Entity the ordering reverses
on the seed-mean convention (CSLS $0.8882$ against $0.8837$; $0.9996$ against
$0.9995$; seed-mean Entity scores are tied), so CRD is not described as uniformly
better. Across the $24$ cells, $21$ are significant after Holm; four of these are the
monotone GMM control, whose AUROC is identical to cosine by construction, so the
independent count is $17$ of $20$. The exceptions are Entity against CSLS, Quantity
against CSLS and Action against TRACED. The Entity gain
over cosine, although significant, is $0.04$ points where every method exceeds
$0.9985$ and is not read as a practical gain.

The comparison against the debiasing baselines is consistent. As a separate family of
ten tests, CRD is better than mean-centred cosine, on the seed-mean score, by
$3.97$ (Attribute), $2.67$ (Quantity), $2.55$ (Action), $1.41$ (Relation) and $0.06$
(Entity) pp, and better than ABTT by $4.22$, $3.56$, $7.51$, $0.66$ and $0.10$\,pp
on the same facets. The Relation-vs-ABTT cell is the weakest ($+0.66$\,pp, CI
$[-0.03,1.32]$) and is not treated as decisive.
Neither baseline exceeds raw cosine on Attribute ($0.6790$ and $0.6765$ against
$0.6806$), so removing first-order feature structure does not reproduce the effect.

\textbf{From ranking to top-$10\%$ label precision.} Precision here is the fraction
of flagged pairs that are noisy. Of the $10\%$ of Attribute test pairs
flagged as most suspicious, $84.9\%$ are noisy, against $76.8\%$ for cosine and
$77.5\%$ for CSLS (paired $+8.0$\,pp over cosine, arithmetic $+8.1$;
cycle-clustered 95\% CI $[7.1,8.9]$). It is not significant elsewhere
($p=0.245$--$0.37$); Entity is a tie. This is a label-precision proxy, not evidence
of a downstream training benefit.

\textbf{Supervised reference.} The same features and training subset are also given
to probes that consume the pair label: a logistic regression on $[u;v;u\odot v]$ and
a two-layer MLP of CRD's width. They are not competitors under the same budget, but
they bound what a classifier extracts. On Attribute the LR probe reaches $0.7234$,
$0.47$\,pp above CRD (bootstrap CI $[0.28,0.60]$); the MLP reaches $0.7214$. On
Quantity the probe is $1.60$ ahead and the MLP $1.85$; on Action the MLP is $7.62$
ahead on the $166$-cycle test split, where $\lambda$ is selected on $80$ internal
validation cycles. With the $320$-row ($80$-cycle) budget, the Action LR probe falls
to $0.5687$ ($-4.59$) while the MLP remains at $0.6517$ ($+3.71$), so the budget
penalty is model-dependent. On Relation the MLP overfits and falls
$3.93$ below CRD. A probe on $[u;v]$ alone is at chance ($0.5000$), as required by
the reciprocal design: any per-side additive score gives identical clean and noisy
sums within a cycle. The interaction $u\odot v$ reaches $0.6977$: the linearly
extractable signal resides there, and nonlinear per-side heads recover only part of
it ($0.6944$/$0.6956$). Trained on the same $8{,}336$-row budget used to select
$\lambda$, the LR probe falls to $0.6751$ on Attribute ($-4.36$ vs CRD) and $0.7949$
on Quantity ($-8.88$): its earlier advantage used $32{,}000$ rows, about four times
CRD's label budget.

\textbf{Head class.} A closed-form ridge version of the same additive form matches
CRD on Attribute and Entity ($0.7187$/$0.9996$), is $0.9$\,pp higher on Quantity,
and ties Relation, so the gain does not require nonlinear heads. The fitted $g$
reaches $R^2=0.326$ on Attribute, while caption length and numeral count reach
$R^2=0.005$, so it is not a surface shortcut.

\subsection{Cross-Backbone Evaluation}
\label{sec:cross}
\begin{table}[t]
\centering
\small
\setlength{\tabcolsep}{3.5pt}
\caption{CRD on three frozen backbones for Attribute (seed-mean AUROC; $\Delta$ in
points). Heads and $\lambda$ are refitted per cell.}
\label{tab:cross}
\begin{tabular}{llccc}
\toprule
Backbone & Facet & Raw & CRD & $\Delta$ \\
\midrule
CLIP B/16  & Attribute & 0.6806 & 0.7187 & +3.81 \\
CLIP L/14  & Attribute & 0.6779 & 0.7139 & +3.60 \\
SigLIP2    & Attribute & 0.7026 & 0.7509 & +4.83 \\
\bottomrule
\end{tabular}
\end{table}

Table~\ref{tab:cross} reports the three Attribute backbones. CRD improves raw AUROC
by $+3.60$ to $+4.83$ points on these cells; we do not extrapolate to unshown facets.
Scaling within CLIP leaves Attribute essentially unchanged (ViT-L/14, $2.9\times$
parameters, $0.6779$ against $0.6806$).

\subsection{Ablation Studies}
\begin{table}[t]
\centering
\small
\setlength{\tabcolsep}{5pt}
\caption{Ablation on Visual Genome Attribute with CLIP ViT-B/16, mean over three
seeds. All rows share the frozen features, training split and $\lambda$ grid.}
\label{tab:ablation}
\begin{tabular}{lccc}
\toprule
Configuration & AUROC & AUPRC & $\Delta$ (pp) \\
\midrule
CLIP cosine (base)      & 0.6806 & 0.6680 & ---    \\
Image-only $a_\phi(u)$  & 0.6944 & 0.6807 & +1.38  \\
Text-only $b_\psi(v)$   & 0.6956 & 0.6904 & +1.50  \\
\textbf{Additive (CRD)} & \textbf{0.7187} & \textbf{0.7177} & \textbf{+3.81} \\
Cycle-level target      & 0.7125 & 0.7077 & +3.19  \\
Shuffled placebo        & 0.6806 & 0.6680 & 0.00   \\
\bottomrule
\end{tabular}
\end{table}

Table~\ref{tab:ablation} isolates the design choices. Image-only reaches $0.6944$ and
text-only $0.6956$, against $0.7187$ for the additive form; the sum of the individual
gains ($1.38+1.50=2.88$) is smaller than the joint gain ($3.81$) by $0.93$, which is
consistent with complementary rather than duplicated information. Supervising on cycle means costs $0.62$ points
($0.7125$), consistent with a target that is constant within the ranked group. The
placebo trains structurally identical heads on permuted cycle means and returns the
baseline ($0.6806$, gain $0.00$) because validation selects $\lambda=0$. Paired
bootstraps give additive minus image-only $+2.43$\,pp, minus text-only $+2.31$,
minus cycle-level $+0.62$ and minus placebo $+3.81$, with all seeds agreeing in sign
and a spread of at most $0.15$\,pp.

\begin{table}[t]
\centering
\footnotesize
\setlength{\tabcolsep}{4pt}
\caption{Same cached features, 1,500-step row target, $\lambda$ grid and three seeds;
$65{,}795$ vs $65{,}665$ parameters. Action excluded (internal split). The additive
column is a same-cache re-run and differs from Table~\ref{tab:main} by at most
$0.003$ AUROC ($0.3$ points). Joint moves Relation 2AFC from $0.6501$ to
$0.6535$; additivity is sufficient, not a ceiling.}
\label{tab:joint}
\begin{tabular}{lcccc}
\toprule
Head & Attribute & Quantity & Entity & Relation \\
\midrule
Additive (CRD) & 0.7187 & 0.8866 & 0.9995 & 0.5538 \\
Joint $[u;v]$  & 0.7182 & 0.8759 & 0.9991 & 0.5537 \\
$\Delta$       & $-0.05$ & $-1.07$ & $-0.04$ & $-0.01$ \\
\bottomrule
\end{tabular}
\end{table}

\subsection{What the Calibration Changes}
\label{sec:analysis}
The 2AFC used here is a cycle-level group comparison: for each cycle it compares the
sum of the two clean scores with the sum of the two noisy scores. Within a cycle both
sides involve the same two images and captions, so summing \eqref{eq:calib} gives
\begin{equation}
\begin{aligned}
\Delta' &= \Delta
- \lambda\Big[\textstyle\sum_{\text{clean}}(a+b+b_0)
- \sum_{\text{noisy}}(a+b+b_0)\Big],
\end{aligned}
\end{equation}
where the bracketed term is identically zero, since both sums contain
$a(u_A)+a(u_B)+b(v_A)+b(v_B)+2b_0$. Any additive correction therefore leaves this
metric unchanged, for any $a$, $b$ and $\lambda$. CRD's 2AFC equals cosine's by
construction ($0.8282$, $11{,}827/14{,}281$ cycles, on Attribute; $0.9649$ Quantity;
$1.0000$ Entity, saturated; $0.7831$, $130/166$, Action; $0.6501$ Relation); CSLS is
invariant for the same reason.

The invariance is a property of the metric, not evidence of harmlessness; it does
not extend to same-image candidate comparison, where only the fixed-side head cancels.
Mean-centred and ABTT cosine are not separable additive corrections and move the
metric (up to $0.0089$ Attribute, $0.1686$ Action), and a monotone rescaling moves it
by $-1.24$ and $-3.51$ points; both metrics are reported.

\subsection{External Transfer}
A head fitted on COCO Entity is applied, without target labels or statistics, to five
protocols. CRD improves on SugarCrepe (+5.08), Flickr noun (+1.04), FOIL (+0.95) and
Flickr random (+0.18), and loses $0.35$ on MM-IMDb; with the VG Attribute source only
SugarCrepe is significantly positive ($+0.98$). Transfer is source-dependent; a
supervised probe is never significantly better.

\subsection{Limitations}
Equation~\eqref{eq:decomp} asserts a predictable additive component; it does not
claim that the component is identifiable as image or text bias. Labels are silver,
with a known substitution failure mode, and are not human-adjudicated. Relation
remains weak (best AUROC $0.5530$, $+1.17$ points). CSLS remains ahead on Quantity
and Entity, and Action is statistically indistinguishable from TRACED. A supervised
classifier is stronger on three facets but needs full labels; CRD is a light-budget
calibration.
